\chapter{Zero-Trust Data Platforms: Identity, Networking, and Secrets}
\index{zero trust}\index{Microsoft Entra ID}\index{managed identity}\index{Private Endpoint}\index{Azure Key Vault}\index{customer-managed keys}%
\begin{objectives}
\begin{itemize}
    \item Apply zero-trust principles---verify explicitly, least privilege, assume breach---to data platform architecture.
    \item Distinguish Microsoft Entra ID identities: users, service principals, and system- versus user-assigned managed identities.
    \item Deploy Private Endpoints and disable public network access so data services exist only inside your VNet.
    \item Contrast Private Endpoints with service endpoints and explain the DNS requirements that make private access actually work.
    \item Centralize secrets in Azure Key Vault with rotation policies and Event Grid-driven expiry notifications.
    \item Evaluate customer-managed keys versus Microsoft-managed keys for regulatory and key-custody requirements.
    \item Design a no-exfiltration data perimeter for a healthcare provider using Synapse data exfiltration protection.
\end{itemize}
\end{objectives}

\begin{crosscloud}
Zero-trust networking and identity-first access are cloud-universal; the product names differ but the architecture---private connectivity, no long-lived secrets, centralized keys---is identical.

\vspace{2mm}
{\footnotesize
\begin{tabular}{@{}p{3.0cm}p{2.9cm}p{3.1cm}p{3.2cm}@{}}
\toprule
\textbf{Capability} & \textbf{AWS} & \textbf{Azure} & \textbf{GCP} \\
\midrule
Identity plane & IAM & Microsoft Entra ID + Azure RBAC & Cloud IAM \\
Workload identity & IAM roles for EC2/EKS & Managed identities & Service account impersonation \\
Private service access & PrivateLink / VPC endpoints & Private Endpoints & Private Service Connect \\
Secrets manager & Secrets Manager & Azure Key Vault & Secret Manager \\
Key custody (CMK) & KMS & Key Vault / Managed HSM & Cloud KMS / HSM \\
Bastion access & Systems Manager Session Manager & Azure Bastion & IAP TCP forwarding \\
\addlinespace
\multicolumn{4}{@{}l@{}}{\textbf{Fabric}: Entra ID + OneLake security + private links; \textbf{Snowflake}: private connectivity (PrivateLink) + Tri-Secret; \textbf{MongoDB}: Atlas private endpoints + customer key management.} \\
\bottomrule
\end{tabular}}

\vspace{1mm}
The exam and interview favorite---``why is a managed identity better than a service principal secret?''---has the same answer on every cloud: no secret material exists to leak, rotate, or commit to git.
\end{crosscloud}

\begin{keyterms}
\textbf{Zero trust}: a security model that verifies every request explicitly, grants least privilege, and assumes the network is hostile.\quad
\textbf{Managed identity}: an Entra ID identity whose lifecycle Azure manages---no credentials in code or configuration.\quad
\textbf{Private Endpoint}: a network interface that binds a PaaS resource to a private IP inside your VNet.\quad
\textbf{Service endpoint}: a VNet-optimized route to a PaaS service that still targets the public IP.\quad
\textbf{CMK}: customer-managed keys---encryption keys you own in Key Vault, with your rotation and revocation policy.
\end{keyterms}

\section{Week 21 Overview: Part VI Opens with the Perimeter}
Month 6 turns the platform from functional to defensible. The security model of the previous twenty weeks---managed identities on linked services, RBAC on storage, Key Vault for secrets---now becomes the explicit subject. Zero trust is the organizing principle: \emph{no implicit trust from network location}; every request authenticated, authorized, and encrypted; the blast radius of any single compromise minimized by least privilege \cite{microsoft2025entra,microsoft2025rbac}.

For data platforms, zero trust resolves to four concrete engineering programs: identity (everything authenticates as a managed identity), networking (nothing is publicly reachable), secrets (nothing long-lived exists outside Key Vault), and custody (keys you can rotate and revoke). This chapter builds all four; Chapter 23 adds the data-privacy layer on top.

\section{Monday: Entra ID, Service Principals, and Managed Identities}
\index{Microsoft Entra ID}\index{service principal}\index{managed identity}

\subsection{The Identity Spectrum}
A \textbf{service principal} is an application identity in Entra ID, authenticated by a client secret or certificate that \emph{you} create, store, rotate, and eventually leak into a log file. A \textbf{managed identity} is a service principal whose credential lifecycle Azure owns: the resource (a VM, a Function, a Data Factory) simply \emph{has} an identity, requests tokens from a local endpoint, and never touches secret material. \textbf{System-assigned} identities live and die with one resource; \textbf{user-assigned} identities are standalone objects shared across resources (useful when many Functions need one identity, or when RBAC must exist before the resource does).

\subsection{The Access Pattern}
The zero-trust access pattern for any service-to-service call: caller has a managed identity $\rightarrow$ identity holds an RBAC role on the target at the tightest scope $\rightarrow$ caller requests a token $\rightarrow$ target validates. No connection strings, no keys, no rotation calendar. Chapters 5, 10, 11, and 18 already used this pattern for COPY loads, linked services, and AML datastores; this section names the principle they were applying.

\begin{pitfall}
\textbf{Broad roles to silence 403s.} Granting \texttt{Owner} or account-wide \texttt{Contributor} because a pipeline ``couldn't read storage'' converts a scoped authorization question into a standing breach. The correct response to a 403 is to identify the missing data-plane action---often \texttt{Storage Blob Data Reader} on one container---not to widen the role.
\end{pitfall}

\section{Tuesday: Private Networking---Endpoints, DNS, and Bastion}
\index{Private Endpoint}\index{service endpoint}\index{Azure Bastion}

\subsection{Private Endpoint vs.\ Service Endpoint}
A \textbf{service endpoint} extends your VNet's identity to a PaaS service over the Azure backbone---the service still has its public IP; traffic is merely routed optimally, and the service firewall must allow your subnet. A \textbf{Private Endpoint} puts a private IP \emph{inside} your VNet that maps to one specific resource; disable public network access on that resource and it ceases to exist on the internet \cite{microsoft2025privateLink}. Zero trust wants the latter.

\subsection{DNS: The Step Everyone Forgets}
A private endpoint without DNS is a trap: clients resolve \texttt{stadls.dfs.core.windows.net} to the \emph{public} IP, traffic fails or leaks, and the misconfiguration is invisible in the portal. The complete pattern is private endpoint + \textbf{private DNS zone} (\texttt{privatelink.dfs.core.windows.net}) + VNet link + the zone group registration on the endpoint. Chapter 24's Terraform listing encodes all four; hand-built setups routinely omit the third or fourth.

\subsection{Administration Without Exposure}
With data services private, administrator access needs its own path: \textbf{Azure Bastion} provides browser-based RDP/SSH to VMs without public IPs on the VMs or NSG openings. The admin plane becomes: Bastion $\rightarrow$ jump VM $\rightarrow$ private endpoints. Nothing in the estate answers a public SYN packet.

\begin{tipbox}
Prove the perimeter from \emph{outside}: after disabling public access, attempt the data-plane call from your workstation and confirm it fails, then from the jump VM and confirm it succeeds. A private endpoint nobody tested is a hope, not a control.
\end{tipbox}

\section{Wednesday: Key Vault, Rotation, and CMK}
\index{Azure Key Vault}\index{secret rotation}\index{customer-managed keys}

\subsection{Secrets, Keys, Certificates}
Azure Key Vault is the custody layer for three object classes: \textbf{secrets} (connection strings, API keys), \textbf{keys} (RSA/EC keys for encryption, including CMKs), and \textbf{certificates} (with managed issuance and renewal) \cite{microsoft2025keyVault}. The zero-trust rule: applications reference secrets \emph{by URI}; the vault enforces access via RBAC; managed identities do the authenticating. A rotated secret must require no code change---consumers read the current version at startup or per use.

\subsection{Automated Rotation}
Key Vault rotation policies plus Event Grid's \texttt{SecretNearExpiry} event close the loop: when a secret nears expiry, an event fires; a Function or Logic App regenerates the credential at the source (for example, a new storage key), writes the new value as a new secret version, and consumers pick it up on next read. The syllabus of Chapter 12's Logic Apps and Chapter 11's Event Grid converge here: rotation is just another event-driven pipeline, one whose payload happens to be a credential.

\subsection{Customer-Managed Keys}
At-rest encryption is on by default with Microsoft-managed keys. \textbf{CMK} moves key custody to your Key Vault: you control rotation, can revoke access by disabling the key (an instant cryptographic kill switch), and can produce the custody evidence regulated industries require. The operational price is real: key availability becomes your availability dependency, and a deleted key without purge protection is a data-loss event. CMK is a compliance answer, not a default posture.

\begin{pitfall}
\textbf{Key Vault references without purge protection.} Disabling purge protection saves nothing and turns an accidental key deletion into permanent ciphertext. Enable soft-delete \emph{and} purge protection on any vault holding CMKs.
\end{pitfall}

\section{Thursday Hands-on Project: A Private-Only Storage Account}
\index{hands-on lab}\index{Private Endpoint!project}
Create an ADLS Gen2 account, bind it to a private endpoint, disable public access, and prove the access pattern from a VM inside the VNet.

\subsection{Build the Perimeter}
\begin{lstlisting}[language=bash,caption={Private endpoint with DNS zone group, then public access disabled.},label={lst:ch21-pep}]
$rg = "rg-de-azure-book-week21"
az group create --name $rg --location eastus
az network vnet create --name vnet-data --resource-group $rg `
  --address-prefix 10.20.0.0/16 --subnet-name snet-data --subnet-prefix 10.20.1.0/24

$adlsId = az storage account show --name <adls> --resource-group $rg --query id -o tsv

az network private-endpoint create --name pep-adls-dfs --resource-group $rg `
  --vnet-name vnet-data --subnet snet-data `
  --private-connection-resource-id $adlsId `
  --group-id dfs --connection-name adls-dfs-conn

az network private-dns zone create --resource-group $rg `
  --name "privatelink.dfs.core.windows.net"
az network private-dns link vnet create --resource-group $rg `
  --zone-name "privatelink.dfs.core.windows.net" --name link-data `
  --virtual-network vnet-data --registration-enabled false

az network private-endpoint dns-zone-group create --resource-group $rg `
  --endpoint-name pep-adls-dfs --name dfszone `
  --private-dns-zone "privatelink.dfs.core.windows.net" --zone-name dfs

az storage account update --name <adls> --resource-group $rg `
  --public-network-access Disabled --default-action Deny
\end{lstlisting}

\subsection{Prove It}
Deploy a small VM into \texttt{snet-data} with Azure Bastion for access. From your workstation: \texttt{az storage fs list} against the account fails (no public path). From the VM via Bastion: the same call succeeds over the private IP, and \texttt{nslookup <adls>.dfs.core.windows.net} returns a 10.20.x.x address---the DNS proof that makes the perimeter real.

\section{Friday Use Case: A No-Exfiltration Perimeter for Healthcare}
\index{use case}\index{data exfiltration protection}
A healthcare provider's compliance requirement: PHI processed in Synapse must not be copyable to any storage outside the governed perimeter, even by a credentialed insider.

\subsection{The Design}
Every data service (ADLS, Synapse, Key Vault, SQL) sits behind private endpoints with public access disabled. The Synapse workspace enables \textbf{data exfiltration protection}: the workspace's managed VNet blocks outbound connections to arbitrary endpoints, so a notebook cannot \texttt{pip install} from a personal repository or POST PHI to an external URL. Egress that \emph{is} required flows through an explicit firewall allowlist, logged. CMK with your rotation evidence satisfies key-custody requirements; Bastion mediates administration; Defender for Cloud (Chapter 23) watches the control plane. The insider threat is bounded: a credentialed user can read what their role allows, but the data has no network path out.

\begin{pitfall}
\textbf{Perimeter without egress control.} Private endpoints stop inbound exposure; they say nothing about a notebook exfiltrating to attacker-controlled storage. Exfiltration protection and egress allowlisting are the outbound half of the same control---deploy both or the perimeter is a one-way valve.
\end{pitfall}

\section{Architectural Decision Record Template}
\index{architectural decision record}
Per platform record: identity model per service (which managed identities, which RBAC roles, at what scopes), the private endpoint inventory with DNS evidence, public-access posture per resource, Key Vault topology (one vault per environment, RBAC mode), CMK scope and rotation ownership, and the egress policy with its allowlist owner.

\section{Operational Deep Dive: Auditing the Perimeter}
\index{audit!security}
The perimeter drifts unless audited: a monthly automated check that lists every data resource with \texttt{publicNetworkAccess} not Disabled, every storage account with shared-key authorization enabled, every Key Vault without purge protection, and every role assignment broader than the approved set. Azure Policy can enforce most of these at deploy time (Chapter 23); the audit catches what predates the policy. Feed both into the same report so posture is one page, not two systems.

\section{Troubleshooting Patterns}
\index{troubleshooting!networking}
Service unreachable after private endpoint creation $\rightarrow$ DNS zone/link/group incomplete; check resolution first. Intermittent private access $\rightarrow$ two DNS paths disagreeing (custom DNS forwarders); trace from the failing client. Pipeline fails after public access disabled $\rightarrow$ the Azure IR is outside the perimeter; use a SHIR inside the VNet (Chapter 10) or a managed VNet IR. Key Vault reference failing after rotation $\rightarrow$ consumer pinned a secret \emph{version} instead of the rolling URI.

\section{Week 21 Lab Rubric}
\index{rubric}
\begin{table}[H]
\centering
\small
\begin{tabular}{@{}p{3.2cm}p{5.0cm}p{5.0cm}@{}}
\toprule
\textbf{Criterion} & \textbf{Meets expectation} & \textbf{Needs improvement} \\
\midrule
Identity & Managed identities everywhere; scoped RBAC & Service-principal secrets; broad roles \\
Networking & Private endpoints + DNS proven from inside and outside & Endpoint without DNS; public access left enabled \\
Secrets & Key Vault references; rotation via events; purge protection & Secrets in variables; no rotation story \\
Verification & Outside-fails/inside-succeeds evidence captured & Controls asserted but never tested \\
\bottomrule
\end{tabular}
\caption{Week 21 rubric for the zero-trust deliverables.}
\label{tab:ch21-rubric}
\end{table}

\section{Interview Q\&A}
\subsection{Concepts and Trade-offs}
\begin{enumerate}
    \item \textbf{Q: Why prefer managed identities over service principal client secrets?}\\
    A: Azure owns the credential lifecycle---no secret to store, rotate, or leak; tokens are issued to the resource on demand.
    \item \textbf{Q: Private Endpoint vs.\ Service Endpoint---what is the key difference?}\\
    A: A private endpoint gives the resource a private IP in your VNet and enables disabling public access; a service endpoint only optimizes routing to the public IP.
    \item \textbf{Q: What does the Key Vault SecretNearExpiry event enable?}\\
    A: Automated rotation: an event-driven pipeline regenerates the credential at the source and publishes a new secret version with no code change.
    \item \textbf{Q: What does VNet data-exfiltration protection in Synapse guard against?}\\
    A: Outbound paths---notebooks or jobs copying data to endpoints outside the governed perimeter, even by credentialed users.
\end{enumerate}

\section{Further Reading}
The Entra ID \cite{microsoft2025entra}, Azure RBAC \cite{microsoft2025rbac}, Private Link \cite{microsoft2025privateLink}, and Key Vault \cite{microsoft2025keyVault} documentations are the primary sources. The Well-Architected security pillar \cite{microsoft2025wellarchitected} frames the zero-trust review; Chapter 23 adds policy enforcement and the privacy layer.

\section*{Key Takeaways}
\begin{itemize}
    \item Zero trust for data platforms: managed identities, private-only networking, vault-custodied secrets, owned keys.
    \item A 403 is a scoping question; broad roles are a standing breach, not a fix.
    \item A private endpoint without its private DNS zone is a trap---test name resolution, not just connectivity.
    \item Disable public access \emph{and} control egress; perimeters are two-way.
    \item Rotation must be event-driven and code-invisible; \texttt{SecretNearExpiry} is the trigger.
    \item CMK buys custody and revocation at the price of availability responsibility; protect the vault with purge protection.
\end{itemize}

\section{Exercises}
\begin{enumerate}
    \item \textbf{(Conceptual)} Explain the token flow when an ADF managed identity reads ADLS Gen2.
    \item \textbf{(Conceptual)} Why does disabling public network access require DNS changes, not just firewall changes?
    \item \textbf{(Conceptual)} When is a user-assigned managed identity preferable to system-assigned?
    \item \textbf{(Hands-on)} Complete the Thursday project, including the outside-fails/inside-succeeds proof.
    \item \textbf{(Hands-on)} Implement a rotation pipeline for a storage key using \texttt{SecretNearExpiry} and a Function.
    \item \textbf{(Hands-on)} Write the monthly posture-audit script and run it against your lab estate.
    \item \textbf{(Design)} Design the egress allowlist for the healthcare Synapse workspace: what legitimately needs outbound?
    \item \textbf{(Design)} Write the ADR for CMK versus Microsoft-managed keys for the PHI estate.
    \item \textbf{(Design)} Plan the migration of a service-principal-based estate to managed identities with zero downtime.
    \item \textbf{(Interview)} Prepare a two-minute answer to ``how would you exfiltrate data from your own platform---and what stops you?''
\end{enumerate}

\section{Capstone Project: Private-Only Lakehouse Perimeter}\label{sec:ch21-capstone}
\index{capstone project}\index{zero trust!capstone}

\begin{capstonebox}
\textbf{Mission:} Convert the lab lakehouse into a zero-trust perimeter: private endpoints and DNS for every data service, public access disabled, managed identities only, Key Vault custody with rotation, Bastion-mediated administration, and an audited posture report.

\textbf{Estimated time:} 4--6 hours.

\textbf{What you will practice:}
\begin{itemize}
    \item Private endpoint + DNS deployment done completely and verifiably.
    \item Identity-first access with scoped RBAC.
    \item Rotation automation and CMK custody.
    \item Posture auditing as a repeatable artifact.
\end{itemize}
\end{capstonebox}

\subsection*{Scenario}
The healthcare provider's security review blocks go-live until the analytics estate is provably unreachable from the internet and provably incapable of unsanctioned egress. Your Part II lakehouse is the conversion target.

\subsection*{Requirements}
\begin{itemize}
    \item \textbf{Network:} private endpoints + DNS zones for ADLS, Synapse, Key Vault; public access disabled; outside-fails/inside-succeeds proof.
    \item \textbf{Identity:} every service-to-service call on managed identities; no shared-key or SAS paths enabled.
    \item \textbf{Secrets:} all secrets in Key Vault by URI reference; one rotation demonstrated via events.
    \item \textbf{Egress:} Synapse exfiltration protection enabled; the allowlist documented.
    \item \textbf{Audit:} the posture script output committed as evidence.
\end{itemize}

\subsection*{Reference Architecture}
\begin{figure}[H]
    \centering
    \resizebox{\textwidth}{!}{%
    \begin{tikzpicture}[node distance=1.2cm]
        \node[stage] (admin) {Admin via\\Azure Bastion};
        \node[stage, right=1.5cm of admin] (vm) {Jump VM\\(VNet)};
        \node[stage, right=1.5cm of vm] (pep) {Private Endpoints\\+ privatelink DNS};
        \node[store, right=1.4cm of pep] (svc) {ADLS / Synapse\\/ Key Vault};
        \node[doc, below=0.9cm of pep] (egress) {Egress allowlist\\+ exfil protection};
        \node[doc, below=0.9cm of svc] (audit) {Posture audit\\(monthly script)};
        \node[stage, left=1.2cm of admin] (inet) {Public internet\\\textbf{denied}};
        \draw[arrow] (admin) -- (vm);
        \draw[arrow] (vm) -- (pep);
        \draw[arrow] (pep) -- (svc);
        \draw[arrow, dashed] (egress) -- (pep);
        \draw[arrow, dashed] (audit) -- (svc);
        \draw[biarrow, draw=red!60!black] (inet) -- node[above,font=\scriptsize\color{red!60!black}]{no path} (admin);
    \end{tikzpicture}%
    }
    \caption{Capstone perimeter: no public path in, allowlisted egress out, every hop authenticated by identity.}
    \label{fig:ch21-capstone-arch}
\end{figure}

\subsection*{Implementation Steps}
\begin{enumerate}
    \item Deploy endpoints and DNS for each service (Listing~\ref{lst:ch21-pep}).
    \item Disable public access and shared-key authorization; retest all pipelines over managed identities.
    \item Configure the vault, rotation policy, and one Event Grid rotation demo.
    \item Enable exfiltration protection; document the allowlist.
    \item Run the posture audit; fix what it finds; commit the clean report.
\end{enumerate}

\subsection*{Deliverables}
\begin{itemize}
    \item CLI/Terraform artifacts for the perimeter and the DNS evidence.
    \item The access proofs and the rotation demo recording.
    \item The posture report and the ADR set for identity, network, secrets, and custody.
\end{itemize}

\subsection*{Acceptance Criteria}
\begin{itemize}
    \item No data service resolves publicly; every in-VNet access path works.
    \item A rotated secret propagates to consumers with zero code or config change.
    \item The posture audit reports zero public-access, shared-key, or purge-protection violations.
\end{itemize}

\subsection*{Stretch Goals}
\begin{itemize}
    \item Convert the entire perimeter to Terraform (preview of Chapter 24) and destroy/rebuild it from code.
    \item Add Defender for Cloud recommendations to the posture report.
    \item Simulate an insider exfiltration attempt from a notebook and document which control stops it.
\end{itemize}
